% ============================================================
%  MTH636 -- Game Theory | Lecture 27
%  Mixed & Behavioral Strategies; SPE Theorem and Remarks
% ============================================================
\documentclass[12pt, a4paper]{article}

% ============================================================
%  MTH636 -- Game Theory
%  Notes Preamble — Color-Coded, TikZ-Enhanced
% ============================================================

% --- Core packages ---
\usepackage[T1]{fontenc}
\usepackage[utf8]{inputenc}
\usepackage{lmodern}
\usepackage[margin=2.5cm, top=2.8cm, bottom=2.8cm]{geometry}
\usepackage{microtype}

% --- Math ---
\usepackage{amsmath, amssymb, amsthm, mathtools}
\usepackage{bm}        % bold math
\usepackage{cancel}    % strike-through in math

% --- Color & Boxes ---
\usepackage[dvipsnames, svgnames, x11names, table]{xcolor}
\usepackage[most]{tcolorbox}
\tcbuselibrary{breakable, skins, theorems}

% --- TikZ & Diagrams ---
\usepackage{tikz}
\usepackage{tikz-qtree}
\usetikzlibrary{arrows.meta, positioning, shapes.geometric, fit, backgrounds,
                decorations.pathreplacing, calc, matrix, trees}
\usepackage{pgfplots}
\pgfplotsset{compat=1.18}

% --- Tables ---
\usepackage{booktabs}
\usepackage{multirow}
\usepackage{array}
\usepackage{tabularx}
\usepackage{makecell}

% --- Typography & Layout ---
\usepackage{enumitem}
\usepackage{parskip}
\usepackage{fancyhdr}
\usepackage{titlesec}
\usepackage{titletoc}
\usepackage{hyperref}
\usepackage{graphicx}
\usepackage{float}
\usepackage{caption}

% --- Color Palette ---
\definecolor{defblue}{HTML}{1A5276}      % Definition header
\definecolor{defbluebg}{HTML}{D6EAF8}    % Definition background
\definecolor{thmgreen}{HTML}{1E8449}     % Theorem header
\definecolor{thmbg}{HTML}{D5F5E3}        % Theorem background
\definecolor{exorange}{HTML}{B7770D}     % Example header
\definecolor{exbg}{HTML}{FEF9E7}         % Example background
\definecolor{warnred}{HTML}{922B21}      % Warning/Key Insight header
\definecolor{warnbg}{HTML}{FADBD8}       % Warning background
\definecolor{noticegray}{HTML}{2C3E50}   % Note header
\definecolor{noticebg}{HTML}{EBF5FB}     % Note background
\definecolor{taketawayyellow}{HTML}{7D6608} % Key Takeaway header
\definecolor{takeawaybg}{HTML}{FDFEFE}   % Takeaway background
\definecolor{accentpurple}{HTML}{6C3483} % Accent / proof
\definecolor{lightgray}{HTML}{F2F3F4}    % shading

% ============================================================
%  TCOLORBOX ENVIRONMENTS
% ============================================================

% -- Definition (Blue) --
\newtcolorbox{defbox}[2][]{
  enhanced, breakable,
  colback=defbluebg, colframe=defblue, coltitle=white,
  fonttitle=\bfseries\sffamily,
  title={Definition: #2},
  attach boxed title to top left={yshift=-2mm, xshift=4mm},
  boxed title style={colback=defblue, rounded corners},
  arc=4pt, boxrule=1.2pt,
  left=6pt, right=6pt, top=8pt, bottom=6pt,
  #1
}

% -- Theorem / Proposition (Green) --
\newtcolorbox{thmbox}[2][]{
  enhanced, breakable,
  colback=thmbg, colframe=thmgreen, coltitle=white,
  fonttitle=\bfseries\sffamily,
  title={Theorem: #2},
  attach boxed title to top left={yshift=-2mm, xshift=4mm},
  boxed title style={colback=thmgreen, rounded corners},
  arc=4pt, boxrule=1.2pt,
  left=6pt, right=6pt, top=8pt, bottom=6pt,
  #1
}

% -- Proposition (Green variant) --
\newtcolorbox{propbox}[2][]{
  enhanced, breakable,
  colback=thmbg, colframe=thmgreen!70!black, coltitle=white,
  fonttitle=\bfseries\sffamily,
  title={Proposition: #2},
  attach boxed title to top left={yshift=-2mm, xshift=4mm},
  boxed title style={colback=thmgreen!70!black, rounded corners},
  arc=4pt, boxrule=1.2pt,
  left=6pt, right=6pt, top=8pt, bottom=6pt,
  #1
}

% -- Example (Orange/Yellow) --
\newtcolorbox{exbox}[2][]{
  enhanced, breakable,
  colback=exbg, colframe=exorange, coltitle=white,
  fonttitle=\bfseries\sffamily,
  title={Example: #2},
  attach boxed title to top left={yshift=-2mm, xshift=4mm},
  boxed title style={colback=exorange, rounded corners},
  arc=4pt, boxrule=1.2pt,
  left=6pt, right=6pt, top=8pt, bottom=6pt,
  #1
}

% -- Key Insight / Warning (Red) --
\newtcolorbox{warnbox}[2][]{
  enhanced, breakable,
  colback=warnbg, colframe=warnred, coltitle=white,
  fonttitle=\bfseries\sffamily,
  title={Key Insight: #2},
  attach boxed title to top left={yshift=-2mm, xshift=4mm},
  boxed title style={colback=warnred, rounded corners},
  arc=4pt, boxrule=1.2pt,
  left=6pt, right=6pt, top=8pt, bottom=6pt,
  #1
}

% -- Note (Gray-Blue) --
\newtcolorbox{notebox}[1][]{
  enhanced, breakable,
  colback=noticebg, colframe=noticegray!60, coltitle=white,
  fonttitle=\bfseries\sffamily,
  title={Note},
  attach boxed title to top left={yshift=-2mm, xshift=4mm},
  boxed title style={colback=noticegray!60, rounded corners},
  arc=4pt, boxrule=0.8pt,
  left=6pt, right=6pt, top=8pt, bottom=6pt,
  #1
}

% -- Key Takeaway --
\newtcolorbox{takebox}[1][]{
  enhanced, breakable,
  colback=takeawaybg, colframe=taketawayyellow, coltitle=white,
  fonttitle=\bfseries\sffamily,
  title={Key Takeaway},
  attach boxed title to top left={yshift=-2mm, xshift=4mm},
  boxed title style={colback=taketawayyellow, rounded corners},
  arc=4pt, boxrule=1.2pt,
  left=6pt, right=6pt, top=8pt, bottom=6pt,
  #1
}

% -- Intuition --
\newtcolorbox{intubox}[1][]{
  enhanced, breakable,
  colback=accentpurple!8, colframe=accentpurple!60, coltitle=white,
  fonttitle=\bfseries\sffamily,
  title={Intuition},
  attach boxed title to top left={yshift=-2mm, xshift=4mm},
  boxed title style={colback=accentpurple!80, rounded corners},
  arc=4pt, boxrule=0.8pt,
  left=6pt, right=6pt, top=8pt, bottom=6pt,
  #1
}

% ============================================================
%  PAYOFF MATRIX HELPER
%  Usage: inside tabular, use \payoff{a}{b} for (a,b)
% ============================================================
\newcommand{\payoff}[2]{$(#1,\;#2)$}
\newcommand{\nash}[2]{\underline{$#1$},\underline{$#2$}}  % underline NE payoffs
\newcommand{\Nset}{\mathbb{N}}
\newcommand{\Rset}{\mathbb{R}}
\newcommand{\Sset}{\mathcal{S}}
\newcommand{\Ui}{u_i}

% ============================================================
%  SECTION FORMATTING
% ============================================================
\titleformat{\section}
  {\large\bfseries\sffamily\color{defblue}}
  {\thesection}{1em}{}[\titlerule]

\titleformat{\subsection}
  {\normalsize\bfseries\sffamily\color{thmgreen!80!black}}
  {\thesubsection}{1em}{}

\titleformat{\subsubsection}
  {\normalsize\itshape\sffamily\color{exorange!80!black}}
  {\thesubsubsection}{1em}{}

% ============================================================
%  HEADER / FOOTER
% ============================================================
\pagestyle{fancy}
\fancyhf{}
\fancyhead[L]{\sffamily\small\color{defblue} MTH636 -- Game Theory}
\fancyhead[R]{\sffamily\small\color{defblue} \leftmark}
\fancyfoot[C]{\sffamily\small\thepage}
\renewcommand{\headrulewidth}{0.4pt}

% ============================================================
%  HYPERREF SETUP
% ============================================================
\hypersetup{
  colorlinks=true,
  linkcolor=defblue,
  urlcolor=thmgreen,
  citecolor=accentpurple,
  bookmarksnumbered=true,
  pdftitle={MTH636 Game Theory Notes},
  pdfauthor={MTH636 Student Notes}
}

% ============================================================
%  MISC HELPERS
% ============================================================
\newcommand{\BR}{\text{BR}}        % Best Response
\newcommand{\NE}{\text{NE}}        % Nash Equilibrium
\newcommand{\SPNE}{\text{SPNE}}    % Subgame Perfect NE
\newcommand{\Dom}{\text{dom}}      % Dominance
\newcommand{\mixs}[1]{\sigma_{#1}} % Mixed strategy
\newcommand{\pures}[1]{s_{#1}}     % Pure strategy


% ============================================================
\begin{document}

% Title
\begin{center}
  {\LARGE\bfseries\sffamily\color{defblue} MTH636 -- Game Theory}\\[6pt]
  {\Large\sffamily Lecture 27: Mixed \& Behavioral Strategies; SPE Theorem}\\[4pt]
\end{center}
\vspace{0.5cm}
\hrule height 1.5pt
\vspace{0.3cm}

% ============================================================
\section{Strategies in Extensive Form Games}

Recall that in an extensive form game $\Gamma$ each player $i$ has a
collection of information sets $H_i$. A \emph{pure strategy} for player $i$
assigns an action to every information set. We now extend this to mixed and
behavioral strategies.

\begin{defbox}{Pure, Mixed, and Behavioral Strategies}
  Let $H_i$ be the set of information sets for player $i$, and let $A(h)$
  denote the set of actions available at information set $h$.
  \begin{itemize}
    \item \textbf{Pure strategy:}
      $s_i : H_i \to \bigcup_{h \in H_i} A(h)$ \quad such that
      $s_i(h) \in A(h)$ for all $h \in H_i$.

    \item \textbf{Behavioral strategy:}
      $b_i : H_i \to \bigcup_{h \in H_i} \Delta A(h)$ \quad such that
      $b_i(h) \in \Delta A(h)$ for all $h \in H_i$. \\
      (A behavioral strategy assigns an \emph{independent} probability
      distribution over actions at \emph{each} information set.)

    \item \textbf{Mixed strategy (over pure strategies):}
      $\sigma_i : S_i \to [0,1]$ \quad such that
      $\sum_{s_i \in S_i} \sigma_i(s_i) = 1$.\\
      (A mixed strategy is a probability distribution over the set of
      \emph{all} pure strategies $S_i$.)
  \end{itemize}
\end{defbox}

\begin{notebox}
  A \textbf{mixed strategy} randomises over complete plans of action (pure
  strategies), while a \textbf{behavioral strategy} randomises independently
  at each information set. In games of \emph{perfect recall}, the two
  concepts are equivalent (Kuhn's Theorem).
\end{notebox}

\subsection*{Example -- Mixed vs.\ Behavioral Strategies}

Consider a game where Player I moves at the root $n_0$ (actions $A$ and $B$)
and then again inside an information set $\{n_1, n_3\}$ (actions $\ell$ and
$r$), and Player II has an information set with actions $L$ and $R$.

\medskip
\noindent Suppose Player I's \emph{mixed strategy} puts weight:
\begin{align*}
  \sigma_1(Al) &= 0.5, & \sigma_1(Ar) &= 0.1,\\
  \sigma_1(Bl) &= 0.3, & \sigma_1(Br) &= 0.1.
\end{align*}
(These are the probabilities on the four pure strategies $Al, Ar, Bl, Br$.)

\medskip
\noindent The corresponding \emph{behavioral strategy} $b_1$ is read off by
conditioning:
\begin{itemize}
  \item At the root: $\Pr[\text{play }A] = 0.5 + 0.1 = 0.6$, \quad
        $\Pr[\text{play }B] = 0.3 + 0.1 = 0.4$.
  \item Inside the information set $\{n_1,n_3\}$:
        $b_1'(\ell) = 0.5$, $b_1'(r) = 0.5$.
\end{itemize}

\medskip
\noindent Given a mixed strategy $\sigma_i$, its \textbf{induced behavioral
strategy} $\bar{\sigma}_i$ is defined as follows.  Let
$\bar{\sigma}_i(h)$ be:
\begin{itemize}
  \item $\sigma_i$ outside $\Gamma_n$ (i.e.\ inherited from the mixed
        strategy in the rest of the game), and
  \item $\bar{\sigma}_i$ inside $\Gamma_n$ (re-normalised conditional
        on reaching $h$).
\end{itemize}

Formally, if $\sigma_i(A) = \sigma_i(A) \cdot \Gamma_n$-weight, then
\[
  \bar{\sigma}_i(A) = \sigma_i(A),\quad
  \bar{\sigma}_i(B) = \sigma_i(B),\quad
  \bar{\sigma}_i(\ell) = \frac{\sigma_i(\ell)}{\sigma_i(\ell)+\sigma_i(r)},\quad
  \bar{\sigma}_i(r) = \frac{\sigma_i(r)}{\sigma_i(\ell)+\sigma_i(r)}.
\]

% ============================================================
\section{Subgame Perfect Equilibrium -- Main Theorem}

\begin{thmbox}{SPE and Reached Subgames (Theorem)}
  Let $\sigma^*$ be a Nash Equilibrium (either in mixed strategies or in
  behavioral strategies) of a game $\Gamma$, and let $\Gamma_n$ be a subgame
  of the game.  If $\rho(n;\, \sigma^*) > 0$ (i.e.\ node $n$ is reached
  with positive probability under $\sigma^*$), then $\sigma^*$ is also a
  Nash Equilibrium for the subgame $\Gamma_n$.
\end{thmbox}

\begin{proof}
  Suppose $\sigma$ restricted to $\Gamma_n$ is \emph{not} an NE of
  $\Gamma_n$. Then there exists a player $i$ and a strategy $\sigma_i'$ in
  the subgame $\Gamma_n$ such that $\sigma_i'$ gives player $i$ strictly more
  utility in $\Gamma_n$.

  Construct a new strategy for player $i$ in the full game:
  \[
    \bar{\sigma}_i =
    \begin{cases}
      \sigma_i  & \text{outside } \Gamma_n, \\
      \sigma_i' & \text{inside } \Gamma_n.
    \end{cases}
  \]

  Since $\rho(n;\, \sigma^*) > 0$, we have
  $\rho\!\bigl(n;\, (\bar{\sigma}_i,\, \sigma_{-i})\bigr) > 0$.
  From here one can show that
  \[
    u_i(\sigma^*) < u_i\!\bigl(\bar{\sigma}_i,\, \sigma_{-i}\bigr),
  \]
  which contradicts the assumption that $\sigma^*$ is an NE of $\Gamma$.
  Therefore $\sigma^*$ restricted to $\Gamma_n$ must be an NE of $\Gamma_n$.

  \medskip
  \noindent (The full proof can be found in Solan's book.)
\end{proof}

% ============================================================
\section{Corollary and Remarks}

\begin{propbox}{Corollary -- Sufficient Condition for SPE}
  Suppose $\sigma$ is an NE of a game $\Gamma$, and under $\sigma$ all
  strategies of all players are given positive probabilities, and all chance
  nodes give positive probability to all edges coming out of them.  Then
  $\sigma$ is an SPE.
\end{propbox}

\begin{notebox}
  The corollary says: if \emph{every} subgame is on the equilibrium path
  (reached with positive probability), then every NE automatically satisfies
  the SPE condition.  This is because the theorem guarantees NE in all
  reached subgames, and under the corollary's conditions all subgames are
  reached.
\end{notebox}

\subsection*{Example Illustrating the Corollary}

Consider the following extensive form game. Player I moves at the root
(actions $A$, $B$, $C$). After $A$, Player II moves (information set
$\{v_1\}$, actions $\ell$ and $r$). After $B$, Player II moves
(information set $\{v_2\}$, actions $\ell$ and $r$). The chance node
$O$ following $C$ assigns probability $\tfrac{1}{4}$ to one branch and
$\tfrac{3}{4}$ to another.

\medskip
\noindent Payoffs from left to right (terminal nodes):
\begin{align*}
  &(4,1),\;(7,2) \quad \text{from } A,\ell \text{ and } A,r,\\
  &(6,3),\;(4,0) \quad \text{from } B,\ell \text{ and } B,r,\\
  &(0,0),\;(8,4),\;(8,8),\;(0,3)\quad\text{from }C\text{ branches.}
\end{align*}

\noindent After resolving the chance node at $C$:
\[
  \frac{1}{4}(8,4) + \frac{3}{4}(8,8) = (8,7).
\]

\noindent The reduced game (after replacing chance node with expected
payoffs) is:
\begin{center}
  \begin{tabular}{c|cc}
           & $\ell$        & $r$           \\
    \hline
    $A$    & \payoff{4}{1} & \payoff{7}{2} \\
    $B$    & \payoff{6}{3} & \payoff{4}{0} \\
    $C$    & \payoff{8}{7} & \payoff{8}{7} \\
  \end{tabular}
\end{center}

\noindent\textbf{Nash Equilibria:} $(C, \ell)$ and $(C, r)$ are both NE.

\noindent\textbf{Backward Induction Equilibria:}
$(C_A,\, \ell R)$ and $(C_A,\, rR)$ are the two backward induction
equilibria of the game.

\begin{thmbox}{Remarks}
  \begin{description}
    \item[R-1.] \textbf{Every backward induction equilibrium is an SPE.}

    \item[R-2.] \textbf{For any extensive form game, there always exists an
          SPE.}\\
          \emph{Proof:} We will discuss this in the next class.
  \end{description}
\end{thmbox}

\begin{intubox}
  Backward induction -- working from the terminal nodes backwards -- always
  produces an SPE. Moreover, the existence of an SPE for \emph{any} finite
  extensive form game (even with imperfect information) is guaranteed.
\end{intubox}

\end{document}
